\section{Quantitative transfer and the four degree profiles}\label{sec:classification}
Fix a uniform axiom template and a stable uniform target family
$G_{b,n}$ as in Section~\ref{sec:systems}. The parameters $k,s,t$ bound
variable, axiom-index and target-index supports, respectively, and $D_G$
bounds the target degree. Put
\[
                 n_* = \max\{n_0,1,kD_G+t\}.
\]
For $\mathrm S\in\{\NS,\PC\}$, the aggregate profile $h_{\mathrm S}(n)$
is the maximum of the least degrees needed to derive the requested
polynomials $G_{b,n}$, $b\in\mathcal B_n$. Zero targets have degree $0$;
an underivable target gives value $\infty$, and the maximum of an empty
request is $0$. We may instead
restrict $\mathcal B$ to one fixed global orbit. At every $n\ge t$ that
orbit has a representative, and equivariance makes its profile the least
degree of any representative. These are two precise choices of a uniform
request. Choosing different global orbits according to $n$ is not covered.

Use the conservative parameters
\begin{equation}\label{eq:profile-parameters}
                      \kappa=2k,\qquad \sigma=s+t.
\end{equation}
The general transfer budget is then $\kappa d+\sigma=2kd+s+t$.
For a fully invariant target, and in particular for refutation, a sharper
version with $\kappa=k$ and $\sigma=s$ is given below. All the profile
deductions are stated in the same notation so that both versions have a
common presentation.

The next theorem is the quantitative lifting statement from which
the classification follows. It fixes the degree budget and varies the
universe size. Its two directions are important: an upward lift keeps
a derivation available on one progression, while the equivalence also
rules out a new derivation appearing later in that residue class.

\begin{theorem}[Periodicity with linear onset]\label{thm:periodicity}
Over any field of prime characteristic $p$, let $h$ be either the NS or
ordinary-PC aggregate profile, or either profile of one fixed target
orbit. For $d\ge1$ put $R(d)=\kappa d+\sigma$. If
\[
 n,n'\ge n_*,\qquad n,n'>2R(d),\qquad
 n\equiv n'\pmod q,\qquad q=p^a>R(d),
\]
then $h(n)\le d$ if and only if $h(n')\le d$. The degree-zero predicate
is constant for $n\ge n_*$.
\end{theorem}
The proof in Section~\ref{sec:target-transfer} transports the actual
coefficient map of the targets and the actual consequence spaces. It does
not require invariant proofs, and it does not infer the PC assertion from
an NS lower bound. The same argument works whether or not the axioms are
consistent.

Define
\begin{equation}\label{eq:NQ}
 \begin{gathered}
 N(0)=n_*,\qquad Q(0)=1,\\
 N(d)=\max\{n_*,2(\kappa d+\sigma)+1\},\qquad
 Q(d)=\min\{p^a:p^a>\kappa d+\sigma\}\quad(d\ge1).
 \end{gathered}
\end{equation}
Here $N(d)$ denotes the onset for a degree threshold; it is not the
larger rank $N=n+q$ used in the deletion construction. $Q(d)$ is the
chosen block size and period. Both are bounded linearly in $d$ beyond
the fixed template threshold.
The onsets are nondecreasing, and $Q(e)$ divides $Q(d)$ for $e\le d$.
Applying Theorem~\ref{thm:periodicity} at every threshold up to $d$ gives
\begin{equation}\label{eq:periodic}
 n,m\ge N(d),\quad n\equiv m\pmod {Q(d)}
 \quad\Longrightarrow\quad
 \min\{h(n),d+1\}=\min\{h(m),d+1\}.
\end{equation}
Here $\min\{\infty,d+1\}=d+1$. For $d\ge1$,
\begin{equation}\label{eq:periodbound}
                         Q(d)\le p(\kappa d+\sigma).
\end{equation}
Indeed, the preceding power of $p$ is at most $\kappa d+\sigma$.
The fixed initial threshold includes the stabilisation of target degrees.
A budget below a target's stabilised degree cannot derive that target;
thus such thresholds are not replaced by tests on a truncated target.

\subsection{The classification theorem}
Define the set of recurring finite degrees by
\[
 E(h)=\{d\in\NN:\text{some }n\ge N(d)\text{ has }h(n)=d\}.
\]
Equation~\eqref{eq:periodic} shows that every member $d$ occurs along an
infinite arithmetic progression:
\[
                  h(n+\ell Q(d))=d\qquad(\ell\ge0).
\]
The degree remains $d$ throughout this progression: lifting does not lower the absolute degree. Its ratio to the universe size tends to zero. Nor does one such progression make the whole profile bounded; other residue classes may remain hard.

Conversely, a value occurring at arbitrarily large sizes occurs beyond
$N(d)$. Thus $E(h)=\rec(h)$, independently of the admissible onset.
The key restriction is that a nonrecurring finite degree $d$ can occur
only before its own onset $N(d)$.

To include the small ranks in an exact formula, write $x_+=\max\{x,0\}$
and put
\[
             \beta(n)=\left\lfloor
                       \frac{(n-2\sigma-1)_+}{2\kappa}
                       \right\rfloor.
\]
The positive part agrees with natural-number subtraction and is irrelevant
once $n\ge2\sigma+1$.

\begin{theorem}[Four-way classification]\label{thm:classification}
For each profile considered in Theorem~\ref{thm:periodicity}, exactly one
of the following four alternatives holds.
\begin{enumerate}[label=\textup{T\arabic*.}]
\item The profile is bounded by a finite constant for all sufficiently
large $n$.
\item $E(h)$ is empty. For every $n\ge n_*$,
\begin{equation}\label{eq:linear}
                              h(n)>\beta(n),
\end{equation}
including $h(n)=\infty$ in the extended order. In particular the finite
values have an $\Omega(n)$ lower bound.
\item $E(h)$ is finite and nonempty, and $h$ is not eventually bounded.
Put $D=\max E(h)$. There is a nonempty proper set
$U\subseteq\{0,\ldots,Q(D)-1\}$ such that, for each $c\in U$, one
finite value $v_c\le D$ satisfies
\[
        h(n)=v_c\quad\text{whenever }n\ge N(D),\quad n\bmod Q(D)=c.
\]
For $n\ge N(D)$ with $n\bmod Q(D)\notin U$, the lower bound
\eqref{eq:linear} holds.
\item $E(h)$ is infinite. For every bound $B$ there are $n,d\ge B$,
with $d\in E(h)$ and $h(n)=d<\infty$, such that
\begin{equation}\label{eq:peaks}
 n\le(p+2)(\kappa d+\sigma),\qquad
 h(n)\ge\frac{n}{(p+2)\kappa}-\frac{\sigma}{\kappa}.
\end{equation}
Also $\liminf_{n\to\infty}h(n)=\min E(h)<\infty$.
\end{enumerate}
In T1 and T3 the liminf likewise equals the least recurring finite value;
in T2 it is $\infty$. All statements leave the fixed finite initial
range unconstrained. NS and PC assign each fixed uniform request the
same type, by Theorem~\ref{thm:comparison}.
\end{theorem}
\begin{proof}
For $n\ge n_*$, the inverse threshold satisfies $N(\beta(n))\le n$.
If $\beta(n)=0$ this follows from $N(0)=n_*$. Otherwise
$2\kappa\beta(n)\le n-2\sigma-1$, which gives the assertion from
\eqref{eq:NQ}. Consequently, if $h(n)\le\beta(n)$, its finite value
belongs to $E(h)$: its own onset is at most $N(\beta(n))$.
If $E(h)=\varnothing$, this is impossible and proves T2, including the
small ranks for which $\beta(n)=0$.

Suppose $E(h)$ is finite and nonempty and let $D=\max E(h)$.
By truncated periodicity, on each class modulo $Q(D)$ the profile has
one fixed value at most $D$, or is always greater than $D$, once
$n\ge N(D)$. The classes of the first kind form $U$. Recurrence puts
some value from $E(h)$ beyond $N(D)$, so $U$ is nonempty. If all classes
belong to $U$, the profile is eventually bounded. Otherwise, at a size
off $U$, the inequality $h(n)\le\beta(n)$ would give a recurring value,
hence a value at most $D$, contrary to the definition of $U$.
This proves the complete T3 split whenever the profile is not bounded.

If $E(h)$ is infinite, choose its elements $d$ arbitrarily large so that
$N(d)=2(\kappa d+\sigma)+1$. A recurring class for $d$ has a
representative in the interval
\[
              N(d)\le n<N(d)+Q(d),\qquad h(n)=d.
\]
Using integrality and \eqref{eq:periodbound} gives
$n\le(p+2)(\kappa d+\sigma)$, hence \eqref{eq:peaks}. The chosen $d$
and $n$ can both exceed any prescribed bound. These are finite degree
values, rather than occurrences of $\infty$.

Whenever $E(h)$ is nonempty, let $e=\min E(h)$. Each of the finitely
many values below $e$ occurs only finitely often, so eventually
$h(n)\ge e$. The value $e$ recurs, proving the liminf formula. If $E(h)$
is empty, the same argument for every finite bound shows $h(n)\to\infty$.
Finally, eventual boundedness forces $E(h)$ to be finite and nonempty:
a recurring value cannot exceed an eventual bound, and at least one of
the finitely many bounded values occurs infinitely often. This proves
exhaustiveness and mutual exclusivity of the four cases.
\end{proof}

The finite liminf in T4 is one number, not an assertion that the profile
is eventually constant. Its infinitely many recurring levels persist
arbitrarily far out. The bound on $N(d)$ provides the degree gap; the
bound on $Q(d)$ places finite representatives of large recurring levels
at sizes linear in their degrees. Both estimates are needed for the peak
statement. The easy classes in T3 may have different finite degrees.
Neither the theorem nor the comparison identifies the NS and PC sets of
recurring degrees.

\begin{corollary}[Sharper invariant-target and refutation bounds]\label{cor:invariant-budget}
For one fully $S_n$-invariant target given by a stable coefficient rule,
Theorems~\ref{thm:periodicity} and~\ref{thm:classification} hold with
$\kappa=k$, $\sigma=s$ and
$n_*=\max\{n_0,1,kD_G\}$. For refutation, take the single target $1$;
then $D_G=0$, $n_*=n_0$, and the transfer budget is $kd+s$.
\end{corollary}
\begin{proof}
The target map has the trivial module $K$ as source. Its image in a
consequence-space quotient is zero or a copy of $K$, and deletion fixes
$K$. Lemma~\ref{lem:constant} therefore detects membership without the
doubled ambient support needed for indexed targets. The actual target
matrix is still transported. The unchanged axiom and PC-closure budget
$kd+s$ gives periodicity, and the preceding profile argument applies
with the smaller parameters.
\end{proof}

\begin{proposition}[The fluctuating case has no fixed periodic split]\label{prop:no-periodic-split}
In T4 there is no period $q_0$, set of its residue classes $U$, finite
bound $B$, and constant $c>0$ such that, at every sufficiently large
size, $h(n)\le B$ on $U$ and $h(n)\ge cn$ off $U$.
\end{proposition}
\begin{proof}
Choose a recurring finite value $e>B$. Its sufficiently late occurrences
cannot lie in $U$, and for $n>e/c$ cannot lie off $U$. This contradicts
recurrence. The argument uses unbounded recurring finite levels, not
merely an unbounded profile.
\end{proof}

\subsection{The polynomial--exponential consequence}\label{sec:dense-types}
With the full explicit coefficient lists of
Definition~\ref{def:dense-size}, let $S(n)$ be the maximum of the minimum
sizes of proofs of the requested targets. Use $S(n)=0$ for an empty
request and $S(n)=\infty$ if any target is underivable. A collection of
proofs for every target has minimum total size between $S(n)$ and
$|\mathcal B_n|S(n)$ when the request is nonempty. Since
$|\mathcal B_n|=n^{O(1)}$, these two size conventions give the same
conclusions below. For one fixed orbit, any representative has the same
minimum size by relabelling.

\begin{corollary}[Polynomial--exponential size bounds]\label{cor:dense-types}
For either NS identities or ordinary-PC derivations of a fixed uniform
request, the type in Theorem~\ref{thm:classification} has the following
size consequences.
\begin{enumerate}[label=\textup{T\arabic*.}]
\item $S(n)\le n^C$ for all sufficiently large $n$, for some constant $C$.
\item $S(n)\ge2^{cn}$ for all sufficiently large $n$, for some $c>0$.
\item On the easy residue classes, $S(n)\le n^C$ for all sufficiently
large $n$; on their complement, $S(n)\ge2^{cn}$, for fixed $C,c>0$.
\item $S(n)\le n^C$ at infinitely many sizes for some fixed $C$, and
$2^{cn}\le S(n)<\infty$ at arbitrarily large sizes for some fixed $c>0$.
Moreover, for every $e>0$ there is $b$ such that
$n^e<S(n)\le n^b$ at infinitely many sizes. There is no eventual periodic
partition with one uniformly polynomial part and one part satisfying an
exponential lower bound.
\end{enumerate}
The exponential statements are lower bounds. They give no universal
matching exponential upper bound. Infinite sizes are interpreted in the
extended order; finiteness is asserted explicitly on the T4 peak
subsequence.
\end{corollary}
\begin{proof}
Write $D_A$ for a fixed upper bound on axiom degree. If the variable
basis contains an object with support size $j>0$, relabelling supplies
a distinct object supported on each $j$-subset of $[n]$. Thus
$v(n)\ge\binom nj\ge n$ for sufficiently large $n$, while
$v(n)=n^{O(1)}$. The target and axiom index counts are polynomial as well.
Lemma~\ref{lem:dense-size} gives polynomial-size proofs for every target
whose derivation degree is bounded by a fixed constant, uniformly over
the request.

If the aggregate degree is a finite value $d>D_A$, some target attains
that maximum. Its minimum proof size is at least
$B_{v(n)}(d-D_A)$ in either system: for NS, a degree-$d$ identity must
contain a multiplier of degree at least $d-D_A$; for PC, any proof must
contain a line of degree at least $d$. Such a target is nonzero, so the
empty NS certificate creates no exception. A linear degree lower bound
therefore yields $2^{\Omega(n)}$. If a target is underivable, its infinite
size satisfies the same lower inequality. This proves T1--T3 and the two
subsequences in T4.

For the additional T4 assertion, choose a recurring degree $d$ with
$\ell=d-D_A>e$. At its infinitely many sufficiently large occurrences,
\[
                   S(n)\ge B_{v(n)}(\ell)
                        \ge\frac{n^{\ell}}{\ell!}>n^e,
\]
while the fixed degree $d$ gives $S(n)\le n^b$ for some $b$.
A periodic polynomial/exponential partition would be contradicted by
choosing $e$ above its polynomial exponent: these occurrences are too
large for its polynomial part and too small for its exponential part.

If every variable support is empty, the entire variable basis is finite
and every permutation fixes it. Each axiom-index orbit consequently
gives one rank-independent polynomial; each target-index orbit does too.
There are finitely many such orbits, and all appear at sufficiently
large rank. The actual axiom and target polynomial sets therefore
stabilise. The aggregate degree is then a finite constant, or is
$\infty$, giving respectively T1 or T2 with infinite size. This covers
the degenerate case without a growing variable count.
\end{proof}

The coefficient-list estimate is \emph{at least} exponential in universe
size. For instance, $v(n)=\Theta(n^2)$ and a line of degree $\Theta(n)$
require $2^{\Theta(n\log n)}$ coefficient entries. Over a fixed finite
field, encoding the axiom and target polynomials by full lists has length
$L(n)=n^{O(1)}$, since their degrees and index supports are fixed.
If $L(n)=O(n^{c_0})$ for a fixed $c_0>0$, the hard-instance bound gives
$2^{\Omega(L(n)^{1/c_0})}$. This is a stretched-exponential lower bound in
input length, with an exponent depending on the template.
